\documentclass[14pt,a4paper,oneside]{extbook}
\usepackage[utf8]{vietnam}
\usepackage{amsmath,amssymb,mathrsfs,fancyhdr,enumerate,multirow,makecell}
\usepackage{tkz-euclide,tikz-3dplot,tikz,tkz-tab}
\usetikzlibrary{arrows.meta,calc,intersections,angles,quotes,backgrounds,shapes.geometric,patterns,through,shadows,positioning}
\usepackage{yhmath}
\usepackage[top=0.7cm, bottom=0.7cm, left=1.0cm, right=1.0cm, headheight=16pt, headsep=3mm, footskip=5mm]{geometry}
\usepackage[hidelinks,unicode]{hyperref}
\usepackage{qrcode}
\usepackage[many]{tcolorbox}
\usepackage{fontawesome}

\pagestyle{fancy}
\fancyhf{}
\fancyhead[L]{\footnotesize\bfseries\color{teal!85!black}\faBook~SỔ TAY KIẾN THỨC TOÁN 4}
\fancyhead[R]{\footnotesize\bfseries\color{blue!85!black}Lớp Toán Thầy Lĩnh -- loptoanthaylinh.vn}
\fancyfoot[C]{\footnotesize\bfseries\color{gray!80!black}-- Sổ tay kiến thức: Bài 26. Luyện tập chung phép cộng và phép trừ --}
\renewcommand{\headrulewidth}{0.4pt}

\begin{document}

% --- BANNER TIÊU ĐỀ BÀI HỌC ---
\noindent
\begin{tcolorbox}[
  enhanced,
  colback=teal!8!white,
  colframe=teal!80!black,
  arc=3mm, boxrule=1pt,
  left=3mm, right=3mm, top=2mm, bottom=2mm,
  halign=center
]
  {\Large\bfseries\color{teal!90!black}\faBookmark\ BÀI 26: LUYỆN TẬP CHUNG PHÉP CỘNG VÀ PHÉP TRỪ}\\[1.5pt]
  {\footnotesize\itshape\color{black!75} Rèn luyện 4 kĩ năng: đặt tính, tính thuận tiện, tìm thành phần chưa biết và giải toán tổng - hiệu trên 1 trang A4}
\end{tcolorbox}

\vspace{2mm}

% --- KHỐI 1 & 2: LÝ THUYẾT & KĨ NĂNG ---
\noindent
\begin{minipage}[t]{0.49\textwidth}
  \begin{tcolorbox}[
    enhanced,
    title={\footnotesize\bfseries\color{white}\faListUl\ 1. Bốn dạng bài tập cốt lõi},
    colback=blue!2!white,
    colframe=blue!75!black,
    arc=2.5mm, boxrule=0.8pt, top=1.8mm, bottom=1.8mm, left=2mm, right=2mm
  ]
    \footnotesize
    $\bullet$ \textbf{Dạng 1: Đặt tính rồi tính}\\
    Thẳng cột các hàng, nhớ mang sang hàng liền trước.\\
    $\bullet$ \textbf{Dạng 2: Tính thuận tiện nhất}\\
    Nhóm cặp tạo số tròn chục, tròn trăm, tròn nghìn.\\
    $\bullet$ \textbf{Dạng 3: Tìm thành phần chưa biết}\\
    $x + b = c \Rightarrow x = c - b$; $x - b = c \Rightarrow x = c + b$.\\
    $\bullet$ \textbf{Dạng 4: Toán có lời văn (Tổng -- Hiệu)}.
  \end{tcolorbox}
\end{minipage}%
\hfill
\begin{minipage}[t]{0.49\textwidth}
  \begin{tcolorbox}[
    enhanced,
    title={\footnotesize\bfseries\color{white}\faCheck\ 2. Kĩ thuật thử lại chuẩn xác},
    colback=teal!2!white,
    colframe=teal!75!black,
    arc=2.5mm, boxrule=0.8pt, top=1.8mm, bottom=1.8mm, left=2mm, right=2mm
  ]
    \footnotesize
    $\bullet$ \textbf{Thử lại phép cộng:}\\
    Lấy \textbf{Tổng trừ đi một số hạng}, nếu ra số hạng còn lại là làm đúng.\\
    $\bullet$ \textbf{Thử lại phép trừ:}\\
    Lấy \textbf{Hiệu cộng với số trừ}, nếu ra số bị trừ là làm đúng.\\
    $\bullet$ \textbf{Kiểm tra nhanh chữ số tận cùng}:\\
    Tổng hàng đơn vị phải khớp với kết quả phép tính.
  \end{tcolorbox}
\end{minipage}

\vspace{2mm}

% --- KHỐI 3: SƠ ĐỒ TƯ DUY TIKZ ---
\noindent
\begin{tcolorbox}[
  enhanced,
  title={\footnotesize\bfseries\color{white}\faSitemap\ SƠ ĐỒ TƯ DUY: HỆ THỐNG HÓA KĨ NĂNG PHÉP CỘNG -- PHÉP TRỪ},
  colback=yellow!2!white,
  colframe=orange!85!black,
  arc=2.5mm, boxrule=0.8pt, top=1.5mm, bottom=1.5mm, left=1mm, right=1mm,
  halign=center
]
  \begin{tikzpicture}[
    scale=0.92,
    core/.style={rectangle, rounded corners=2.5mm, draw=red!80!black, fill=red!15, thick, inner sep=4pt, font=\scriptsize\bfseries, align=center},
    branch/.style={rectangle, rounded corners=2mm, draw=blue!75!black, fill=blue!10, thick, inner sep=3pt, font=\scriptsize\bfseries, align=center, text width=3.6cm},
    arr/.style={-{Stealth[scale=0.8]}, thick, draw=black!70}
  ]
    \node[core] (root) at (0,0) {\textbf{PHÉP CỘNG \& TRỪ}\\Toán 4 Chủ đề 5};

    \node[branch] (b1) at (5.5, 1.35) {\textbf{1. Đặt tính rồi tính}\\Căn thẳng cột bên phải\\Cộng/trừ từ PHẢI $\rightarrow$ TRÁI};
    \node[branch, draw=teal!80!black, fill=teal!10] (b2) at (11.5, 1.35) {\textbf{2. Tính thuận tiện}\\Giao hoán + Kết hợp\\Ghép cặp tận cùng tròn chục};

    \node[branch, draw=orange!80!black, fill=orange!12] (b3) at (5.5, -1.35) {\textbf{3. Tìm thành phần}\\Số bị trừ = Hiệu + Số trừ\\Số trừ = Số bị trừ - Hiệu};
    \node[branch, draw=purple!80!black, fill=purple!12] (b4) at (11.5, -1.35) {\textbf{4. Toán Tổng -- Hiệu}\\Số bé = $(\text{T} - \text{H}) : 2$\\Số lớn = $(\text{T} + \text{H}) : 2$};

    \draw[arr] (root.east) -- (b1.west);
    \draw[arr] (root.east) -- (b2.west);
    \draw[arr] (root.east) -- (b3.west);
    \draw[arr] (root.east) -- (b4.west);
  \end{tikzpicture}
\end{tcolorbox}

\vspace{2mm}

% --- KHỐI 4: VÍ DỤ MINH HỌA & BẪY TRẮC NGHIỆM ---
\noindent
\begin{minipage}[t]{0.53\textwidth}
  \begin{tcolorbox}[
    enhanced,
    title={\footnotesize\bfseries\color{white}\faCheckCircle\ Ví dụ thực hành chuẩn mực},
    colback=green!2!white,
    colframe=green!60!black,
    arc=2.5mm, boxrule=0.8pt, top=1.8mm, bottom=1.8mm, left=2.2mm, right=2.2mm
  ]
    \footnotesize
    $\bullet$ \textbf{Ví dụ 1 (Tính thuận tiện):}\\
    $4\,560 + 2\,890 + 5\,440 = (4\,560 + 5\,440) + 2\,890$\\
    $= 10\,000 + 2\,890 = \mathbf{12\,890}$.\\[1mm]
    $\bullet$ \textbf{Ví dụ 2 (Tìm $x$):}\\
    $x - 1\,450 = 3\,550 \Rightarrow x = 3\,550 + 1\,450 = \mathbf{5\,000}$.
  \end{tcolorbox}
\end{minipage}%
\hfill
\begin{minipage}[t]{0.45\textwidth}
  \begin{tcolorbox}[
    enhanced,
    title={\footnotesize\bfseries\color{white}\faExclamationTriangle\ Mẹo tránh bẫy trắc nghiệm},
    colback=red!2!white,
    colframe=red!70!black,
    arc=2.5mm, boxrule=0.8pt, top=1.8mm, bottom=1.8mm, left=2.2mm, right=2.2mm
  ]
    \footnotesize
    $\bullet$ \textbf{Bẫy tìm Số trừ và Số bị trừ:}\\
    --- Muốn tìm Số bị trừ: Lấy \textbf{Hiệu $+$ Số trừ}.\\
    --- Muốn tìm Số trừ: Lấy \textbf{Số bị trừ $-$ Hiệu}.\\
    Nhiều bạn máy móc cứ thấy dấu trừ là làm phép cộng!
  \end{tcolorbox}
\end{minipage}

\end{document}
